\justifying
\NSFCBodyText

\subsubsection{已具备的实验条件}

申请人已确认能够稳定使用NVIDIA A40显卡8张和NVIDIA GeForce RTX 3090 Ti显卡6张，并具备Linux、Python、深度学习框架、Git与LaTeX等软件工具条件。上述资源可支持已有模型推理、任务重放和评价实验，项目无需另行采购大型计算设备。实际模型版本和部署规模在预实验中根据资源占用确定，不以全部算力长期独占为前提。

项目可从具有明确使用条件的公开课程资料和论文资源构建任务，配合本地模拟检索、工具调用和记忆环境形成可观察轨迹。公开可获取不等于本课题已拥有数据使用授权，任务整理时将逐项核对许可、引用与可发布范围。现有条件主要支撑公开资料与可控沙箱研究，真实教学数据另行落实。

\subsubsection{尚缺少的实验条件及拟解决的途径}

本课题尚需落实\textbf{【待补 F-EDU-01：合作人与教学规则标注人员】}及\textbf{【待补 F-EDU-02：课程场景、数据使用授权和伦理审查结论】}。拟在研究启动阶段明确合作分工，邀请具有相关教学经验的人员完善任务规则与标注手册，并在必要审批完成后组织小规模案例验证。涉及真人参与的活动应与课程成绩和正式教学评价分离，参与与退出不影响学生权益。

若上述条件未按计划落实，以公开材料、虚拟角色和模拟工具完成主要方法实验，并将案例明确标注为近真实验证。需调用外部模型时仅使用可公开或获准处理的数据，记录成本；拟申请的5万元是当前材料中的建议额度，实际经费与成果要求以正式任务书为准。

\subsubsection{利用研究基地的计划与落实情况}

本项目拟依托教育大数据应用技术国家工程研究中心开放课题开展合作，合作重点为教育风险定义、教学支持尺度和案例分析。当前材料中的合作方向属于拟议安排，尚不能表述为已签订合作协议、已获得中心专属设备或已接入其教学数据。申请人已有算力支撑核心技术工作，中心合作条件落实后用于完善教育情境解释与案例验证。
